\chapter{Kelangsungan Bisnis dan Ketahanan Operasional}

\section*{Tujuan Pembelajaran}
\begin{itemize}
  \item Melakukan BIA (\textit{Business Impact Analysis}) untuk menentukan layanan kritis beserta RTO dan RPO-nya.
  \item Memilih strategi kelangsungan yang sepadan dengan dampak dan biaya.
  \item Merancang skenario latihan krisis yang menguji keputusan, bukan hanya teknologi.
\end{itemize}

\section{Pendahuluan}

% TODO
% Cakupan: SIMRS mati selama enam jam, antrean pasien mengular, dan tidak ada yang tahu prosedur manual mana yang berlaku.

\section{Ketahanan Operasional sebagai Tanggung Jawab Direksi}

% TODO
% Cakupan: beda manajemen insiden, kelangsungan bisnis, dan pemulihan bencana; layanan kritis di rumah sakit dan kaitannya dengan keselamatan pasien; peran direksi dalam menetapkan toleransi gangguan.

\section{ISO 22301:2019}

% TODO
% Cakupan: struktur sistem manajemen kelangsungan bisnis; Amandemen 1:2024 tentang perubahan iklim; revisi ISO 22301 yang sedang berjalan (periksa ulang); kaitan dengan kontrol kelangsungan pada ISO/IEC 27001.

\section{Analisis Dampak Bisnis}

% TODO
% Cakupan: dampak menurut waktu; MTPD, RTO, RPO, dan MBCO (Minimum Business Continuity Objective); ketergantungan pada sistem, orang, tempat, dan vendor.
% Kasus Medikara: pendaftaran pasien, IGD, farmasi, laboratorium, klaim BPJS Kesehatan, dan telekonsultasi.

\section{Strategi Kelangsungan}

% TODO
% Cakupan: strategi untuk orang, tempat, teknologi, dan pemasok; prosedur manual pengganti; pertimbangan biaya terhadap dampak; disaster recovery teknis dibahas rinci di IT40713.

\section{Insiden Siber dan Manajemen Krisis}

% TODO
% Cakupan: Peraturan BSSN 1/2024 tentang Pengelolaan Insiden Siber dan Peraturan BSSN 2/2024 tentang Manajemen Krisis Siber; tim tanggap insiden; komunikasi krisis kepada pasien, regulator, dan media; kaitan dengan notifikasi menurut UU PDP.

\section{Latihan dan Pengujian}

% TODO
% Cakupan: jenis latihan (tabletop, simulasi, uji penuh); merancang skenario yang menguji keputusan; pelajaran setelah latihan dan tindak lanjutnya.
% Pendalaman teknis: menguji apakah RTO dan RPO dapat dipenuhi rancangan replikasi dan cadangan di IT40713, lalu mencatat selisihnya sebagai risiko.

\section{Aktivitas: BIA dan Strategi Kelangsungan Layanan}

% TODO
% Tujuan aktivitas: menetapkan layanan mana yang harus pulih lebih dulu, seberapa cepat, dan dengan cara apa.
% Peran dalam tim: anggota berlatar kesehatan atau operasional memimpin penilaian dampak; analis teknologi menilai kelayakan RTO dan RPO terhadap arsitektur target; analis risiko menautkan hasil ke risk register; pemilik konteks memvalidasi prosedur manual.
% Langkah: pilih lima sampai tujuh layanan; nilai dampak menurut waktu; tetapkan MTPD, RTO, dan RPO; petakan ketergantungan; pilih strategi untuk dua layanan paling kritis; susun skenario tabletop serangan ransomware pada SIMRS.
% Luaran dan tonggak capstone: BIA, strategi kelangsungan untuk dua layanan kritis, dan skenario latihan; masuk C3.
% Berkas sources yang direncanakan: sources/chapter10/templat-bia.xlsx, sources/chapter10/skenario-tabletop-ransomware.md

\section{Ringkasan}

% TODO

\section{Latihan}

\begin{enumerate}
  \item Jelaskan perbedaan RTO dan RPO dengan contoh layanan farmasi rumah sakit.
  \item Mengapa RTO untuk semua layanan tidak boleh ditetapkan nol? Kaitkan dengan biaya dan \textit{risk appetite}.
  \item Rancang tiga pertanyaan keputusan yang harus dijawab direksi dalam latihan \textit{tabletop} serangan \textit{ransomware}.
\end{enumerate}
